\documentclass{ximera}
\title{The derivative}
\begin{document}
\begin{abstract}
The derivative as a function, notation for the derivative, differentiability, and the fact that differentiable functions are continuous. Textbook §2.8.
\end{abstract}
\maketitle

Recall that we've so far been looking at the derivative of a function $f$ at a fixed number $a$
\[
f^{\prime}(a)=\lim _{h \rightarrow 0} \frac{f(a+h)-f(a)}{h} .
\]

What if we want to look at the derivative at an arbitrary point, not just at $a$? For this we replace the number $a$ by a variable.
\[
f^{\prime}(x)=\lim _{h \rightarrow 0}\answer{\frac{f(x+h)-f(x)}{h}}
\]

We already did this when we were looking at acceleration. Notice that this allows us to look at $f^{\prime}$ as a function! This function, $f^{\prime}$, is called the derivative of $f$.

\begin{example}\label{example:4-14}
Let's do a test run with a super easy example. Find the derivative of $f(x)=x$.
% work: ((x+h)-x)/h = h/h = 1
\[
f^{\prime}(x)=\lim _{h \rightarrow 0}\frac{f(x+h)-f(x)}{h}=\lim _{h \rightarrow 0}\frac{(x+h)-x}{h}=\lim _{h \rightarrow 0}\answer{1}
\]
Therefore $f^{\prime}(x)=1$.
\end{example}

\begin{exercise}\label{exercise:4-15}
Try a more complicated example with a partner. Find the derivative of $f(x)=x^{2}-x$.
% work: ((x+h)^2-(x+h)-(x^2-x))/h = (2xh+h^2-h)/h = 2x+h-1 -> 2x-1
\[
f^{\prime}(x)=\lim _{h \rightarrow 0}\frac{\left((x+h)^{2}-(x+h)\right)-\left(x^{2}-x\right)}{h}=\answer{2x-1}
\]
\end{exercise}

There are a lot of different ways that people write the derivative $f^{\prime}(x)$. This comes from the history of derivation!
\[
f^{\prime}(x)=\dot{y}=\frac{d y}{d x}=\frac{d f}{d x}=D f
\]

\begin{definition}\label{definition:4-16}
We say that a function $f$ is \dfn{differentiable at $a$} if $f^{\prime}(a)$ exists.
\end{definition}

\begin{example}\label{example:4-17}
Where is the following function differentiable?
\[
f(x)= \begin{cases}1 & \text { if } x \geq 0 \\ -1 & \text { if } x<0\end{cases}
\]
% work: for x != 0 the function is constant near x, so f'(x)=0; at x=0, (f(0+h)-f(0))/h = -2/h for h<0, which has no limit
The function $f$ is differentiable \wordChoice{\choice[correct]{everywhere except at $x=0$}\choice{everywhere}\choice{only at $x=0$}\choice{nowhere}}.
\end{example}

Being differntiable, actually tells us whether our function is continuous or not!

\begin{theorem}\label{theorem:4-18}
If $f$ is differentiable at $a$, then $f$ is continuous at $a$.
\end{theorem}
\begin{freeResponse}\end{freeResponse}

We see this already in our example!
\begin{image}
\includegraphics[width=0.6\textwidth]{b48a5324-93dc-45ee-b7c2-dc0b5706bda6-057_883_928_501_597}
\end{image}
\begin{freeResponse}\end{freeResponse}

\begin{warning}
The opposite direction is not true! There are functions that are continuous at a point but not differentiable! There is an example of such a function in the book.
\end{warning}

Intuitively, there are three main cases when a function is not differentiable:
\begin{image}
\includegraphics[width=0.4\textwidth]{b48a5324-93dc-45ee-b7c2-dc0b5706bda6-057_339_460_1633_520}
\end{image}
\begin{image}
\includegraphics[width=0.4\textwidth]{b48a5324-93dc-45ee-b7c2-dc0b5706bda6-057_337_460_1633_1061}
\end{image}
\begin{image}
\includegraphics[width=0.4\textwidth]{b48a5324-93dc-45ee-b7c2-dc0b5706bda6-057_335_455_1972_834}
\end{image}
\begin{freeResponse}\end{freeResponse}

\end{document}
